% ===== 原文 pp.33–34 =====
% tr-05 片段 b：Virasoro 算子的一般形式、定理 14（Witten--Kontsevich 猜想）、
%                定理 15（Virasoro 约束）与广义 string/dilaton 方程（p.33–34）
% 本块无编号公式；编号陈述为定理 14、定理 15。

% [原文 p.33]
在简单的条件下，算子的定义可以唯一地延拓，从而得到 $\{L_n\}$ 子代数的一个表示。
一般的算子是
\[
\begin{aligned}
L_n={}&-\frac{(2n+3)!!}{2^{n+1}}\frac{\partial}{\partial t_{n+1}}
+\sum_{j=0}^{\infty}\frac{(2j+2n+1)!!}{(2j-1)!!\,2^{n+1}}t_j
\frac{\partial}{\partial t_{j+n}}\\
&+\frac{\lambda^2}{2}\sum_{j=0}^{n-1}\frac{(2j+1)!!(2n-2j-1)!!}{2^{n+1}}
\frac{\partial^2}{\partial t_j\partial t_{n-j-1}},
\end{aligned}
\]
其交换子为 $[L_n,L_m]=(n-m)L_{n+m}$。

\begin{thm}{14}
[Mir07b, Theorem 6.1]. Witten--Kontsevich 猜想：Virasoro 约束。对 $n\ge -1$，则
$L_n(e^{\mathbf F})=0$。
\end{thm}

\begin{proof}
对指数多重指标 $\mathbf k=(k_1,\dots,k_n)$，体积递推公式成为系数关系
\[
(2k_1+1)V_{g,n}(L)[\mathbf k]
=\cA^{con}_{g,n}(L)[\mathbf k]+\cA^{dcon}_{g,n}(L)[\mathbf k]+\cB_{g,n}(L)[\mathbf k].
\]
% [Tutor Note: p.33 的 $V_{g,n}(L)[\mathbf k]$ 表示取单项式 $L^{2\mathbf k}=\prod_i L_i^{2k_i}$ 的系数，原文未明说；把体积递推按 $L$ 的总次数分层，最高次（$|\mathbf k|=3g-3+n$）那一层就是这里的“首项系数关系”。右端三项分别对应体积递推的 $\cA^{con}$（切开边界后连通）、$\cA^{dcon}$（切开后不连通）与 $\cB$ 三类贡献。重排后得到的是 $L_{k_1-1}(e^{\mathbf F})=0$：由 $k_1\ge 0$ 得 $k_1-1\ge -1$，恰落在 $L_n$ 的定义范围 $n\ge -1$ 内，这是一处可自检的下标一致性。]

首项系数关系取如下的显式形式（按照边界的标号方式，穿刺点标号为 $1,\dots,n$）
\[
\begin{aligned}
(2k_1+1)!!\langle\tau_{k_1}\cdots\tau_{k_n}\rangle
&=\frac12\sum_{i+j=k_1-2}(2i+1)!!(2j+1)!!
\sum_{I\subset\{2,\dots,n\}}\langle\tau_i\tau_{\mathbf k_I}\rangle
\langle\tau_j\tau_{\mathbf k_{I^c}}\rangle\\
&\quad+\frac12\sum_{i+j=k_1-2}(2i+1)!!(2j+1)!!
\langle\tau_i\tau_j\tau_{k_2}\cdots\tau_{k_n}\rangle\\
&\quad+\sum_{j=2}^{n}\frac{(2k_1+2k_j-1)!!}{(2k_j-1)!!}
\langle\tau_{k_2}\cdots\tau_{k_1+k_j-1}\cdots\tau_{k_n}\rangle.
\end{aligned}
\]
重排该显式关系可得 $L_{k_1-1}(e^{\mathbf F})=0$。
\end{proof}

Mulase 与 Safnuk 考虑 $\kappa_1$ 类与 $\psi$ 类的各种组合的相交数的生成函数 [MS08]
\[
\mathbf G(s,t_0,t_1,\dots)
=\sum_g\big\langle e^{s\kappa_1+\sum_j t_j\tau_j}\big\rangle_g
=\sum_g\sum_{m,\{d_j\}}\langle\kappa_1^m\tau_0^{d_0}\tau_1^{d_1}\cdots\rangle_g
\frac{s^m}{m!}\prod_{j=0}^{\infty}\frac{t_j^{d_j}}{d_j!},
\]
其中同样地，除 $3g-3+n$ 乘积之外的乘积都定义为零。Mulase 与 Safnuk 用体积递推
与项的重排证明了下面的结论。

% [原文 p.34]
\begin{thm}{15}
[MS08, Thrm. 1.1] Virasoro 约束。对每个 $k\ge -1$，定义
\[
\begin{aligned}
\cV_k={}&-\frac12\sum_{i=0}^{\infty}(2(i+k)+3)!!\frac{(-2s)^i}{(2i+1)!}
\frac{\partial}{\partial t_{i+k+1}}\\
&+\frac12\sum_{j=0}^{\infty}\frac{(2(j+k)+1)!!}{(2j-1)!!}t_j
\frac{\partial}{\partial t_{j+k}}\\
&+\frac14\sum_{\substack{d_1+d_2=k-1\\ d_1,d_2\ge 0}}
(2d_1+1)!!(2d_2+1)!!\frac{\partial^2}{\partial t_{d_1}\partial t_{d_2}}
+\frac{\delta_{k,-1}t_0^2}{4}+\frac{\delta_{k,0}}{48},
\end{aligned}
\]
其中双阶乘与 Kronecker delta 函数 $\delta_{*,*}$ 的含义如前。则
\begin{itemize}[leftmargin=1.6em,itemsep=2pt,topsep=2pt]
\item 诸算子 $\cV_k$ 满足 Virasoro 交换关系 $[\cV_n,\cV_m]=(n-m)\cV_{n+m}$；
\item 生成函数 $\mathbf G$ 满足 $\cV_k(e^{\mathbf G})=0$，$k\ge -1$。
\end{itemize}
初值条件与第二组方程唯一确定该生成函数。
\end{thm}

% [Tutor Note: 定理 15 取 $k=-1$ 可作结构自检：此时 $d_1+d_2=-2$ 无解，第三个求和为空，$\cV_{-1}$ 只剩前两个求和加 $\frac{\delta_{-1,-1}t_0^2}{4}=\frac{t_0^2}{4}$，正对应 string 方程所属算子的“$t_0^2$ 项”（比较 p.32 的 $L_{-1}$ 中的 $\frac{\lambda^{-2}}{2}t_0^2$）。另取 $s=0$：第一个求和只剩 $i=0$ 项，系数为 $-\frac{(2k+3)!!}{2}$，与 p.33 的 $L_k$ 首项 $-\frac{(2k+3)!!}{2^{k+1}}$ 相差因子 $2^k$——两套算子 $\cV_k$ 与 $L_k$ 并非逐项相同，比较前须先对齐 $t_j$ 与 $s$ 的归一化；而 $k=-1$ 时第二个求和在 $j=0$ 处出现 $\partial/\partial t_{-1}$，按此类公式的惯例应把负指标的 $t$ 与 $\partial_t$ 取为零。]

为了直接显示 $\kappa_1$ 类与 $\psi$ 类的相交数等同于 $\psi$ 类的相交数，Mulase 与
Safnuk 给出
\[
\mathbf G(s,t_0,t_1,t_2,t_3\dots)
=\mathbf F(t_0,t_1,t_2+\gamma_2,t_3+\gamma_3,\dots),
\]
其中 $\gamma_j=-(-s)^{j-1}/((2j+1)j!)$，[MS08, Thrm. 1.2]。该关系的一个显式证明还来自
Faber 的一个公式，它把 kappa 类用更多穿刺点的模空间上的 psi 类表示出来。

在其学位论文 [Do08] 中，Norman Do 给出一个
\[
\text{广义 string 方程}\qquad
V_{g,n+1}(L,2\pi i)=\sum_{k=1}^{n}\int L_kV_{g,n}(L)\,\dd L_k,
\]
以及
\[
\text{广义 dilaton 方程}\qquad
\frac{\partial V_{g,n+1}}{\partial L_{n+1}}(L,2\pi i)
=2\pi i(2g-2+n)V_{g,n}(L),
\]
其中左端把第 $(n+1)^{st}$ 个边界长度取为 $2\pi i$，且 $L=(L_1,\dots,L_n)$。由定理 4，
第二个方程对适当选取的非负多重指标 $\alpha=(\alpha_1,\alpha_2,\dots,\alpha_n)$ 与整数 $m$，
等价于关系
\[
\int_{\Mgnb}\psi_1^{\alpha_1}\psi_2^{\alpha_2}\cdots\psi_n^{\alpha_n}
\psi_{n+1}(\kappa_1-\psi_{n+1})^m
=(2g-2+n)\int_{\Mgnb}\psi_1^{\alpha_1}\psi_2^{\alpha_2}\cdots\psi_n^{\alpha_n}\kappa_1^m.
\]
广义方程的一个证明基于 psi 与 kappa 类的拉回关系，以及关于类之像的一般考量。
特别地，对
